\section{Toán học, công cụ để làm khoa học}

\ 

Chắc các bạn đọc từng nghe ngạn ngữ Anh: ``Một bức hình thì giá trị bằng cả ngàn câu từ'', và thực tế, chúng ta không thể phủ nhận tính súc tích và dễ hiểu của một tranh minh họa được đặt đúng vị trí. Các phương trình toán học cũng như vậy. Hãy so sánh hai cách biểu diễn, cho cùng một định luật quen thuộc:

\begin{table}[H]
   \centering
   \caption{So sánh cách biểu diễn bằng lời và bằng phương trình toán}
   \begin{tblr}{
      colspec = {Q[m,l,7cm] Q[m,c,7cm]}, % Q là cột thông minh
      % m: căn giữa dọc, l: căn trái ngang, c: căn giữa ngang
      hlines, vlines,
      rows = {1.5ex}, % Độ cao dòng
      hline{2} = {1.5pt}, 
      vline{2} = {1.5pt}
   }
      \centering\textbf{Biểu diễn bằng lời} & \textbf{Biểu diễn bằng phương trình} \\
      Gia tốc của một vật thì tỉ lệ thuận với lực tác động lên vật đó và tỉ lệ nghịch với khối lượng của nó. & 
      $\vec{a} = \frac{\vec{F}}{m}$. \\
      Độ dịch chuyển của vật bằng tổng của tích vận tốc ban đầu với thời gian và một nửa tích của gia tốc với bình phương thời gian. & $\Delta x = v_0 \cdot t + \frac{a \cdot t^2}2$.
   \end{tblr}
\end{table}

Có thể thấy rằng, phương trình toán đạt một sự súc tích, ngắn gọn, và mạch lạc hơn so với viết bằng lời. Có thể bạn đọc sẽ phản biện rằng tác giả đã bỏ qua phần giải thích kí hiệu. Nếu có chúng, phương pháp biểu diễn bằng toán học sẽ không còn ngắn gọn nữa. Chúng ta hãy cùng kiểm chứng điều đó.

\begin{table}[H]
   \centering
   \caption{So sánh cách biểu diễn bằng lời và bằng phương trình toán có giải thích}
   \begin{tblr}{
      colspec = {Q[m,l,7cm] Q[m,c,7cm]},
      hlines, vlines,
      hline{2} = {1.5pt}, 
      vline{2} = {1.5pt}
   }
      \centering\textbf{Biểu diễn bằng lời} & \textbf{Biểu diễn bằng phương trình} \\
      Độ dịch chuyển của vật bằng tổng của tích vận tốc ban đầu với thời gian và một nửa tích của gia tốc với bình phương thời gian. & $\Delta x = v_0 \cdot t + \frac{a \cdot t^2}2$, trong đó \begin{itemize}
         \item $\Delta x$: độ dịch chuyển của vật,
         \item $v_0$: vận tốc đầu của vật,
         \item $t$: thời gian chuyển động,
         \item $a$: độ lớn của gia tốc. 
      \end{itemize}
   \end{tblr}
\end{table}

Có vẻ sự phản biện là đúng, nhưng đó là khi chúng ta chỉ sử dụng kết quả này một lần. Hãy cho thêm một vài tính chất có thể chúng ta sẽ quan tâm.

\begin{table}[H]
   \centering
   \caption{So sánh cách biểu diễn bằng lời và bằng phương trình toán (nhiều tính chất)}
   \begin{tblr}{
      colspec = {Q[m,l,7cm] Q[m,c,7cm]},
      hlines, vlines,
      hline{2} = {1.5pt}, 
      vline{2} = {1.5pt}
   }
      \centering\textbf{Biểu diễn bằng lời} & \textbf{Biểu diễn bằng phương trình} \\
      Độ dịch chuyển của vật bằng tổng của tích vận tốc ban đầu với thời gian và một nửa tích của gia tốc với bình phương thời gian. Khi đó, \begin{itemize}
         \item Vận tốc đầu bằng hiệu của độ dịch chuyển và một nửa tích của gia tốc với bình phương thời gian, sau đó đem chia cho thời gian;
         \item Gia tốc bằng hai lần hiệu giữa độ dịch chuyển và tích của vận tốc đầu với thời gian, sau đó chia cho bình phương thời gian.
         \item Thời gian chuyển động bằng thương số của hiệu giữa căn bậc hai của tổng bình phương vận tốc đầu cộng hai lần tích gia tốc và độ dịch chuyển so với vận tốc đầu, chia cho gia tốc.
      \end{itemize} & $\Delta x = v_0 \cdot t + \frac{a \cdot t^2}2$, trong đó \begin{itemize}
         \item $\Delta x$: độ dịch chuyển của vật,
         \item $v_0$: vận tốc đầu của vật,
         \item $t$: thời gian chuyển động,
         \item $a$: độ lớn của gia tốc. 
      \end{itemize}
      \begin{flushleft}
         Khi đó, $v_0 = \frac{\Delta x - \frac{1}{2}at^2}{t}$, $a = \frac{2(\Delta x - v_0 t)}{t^2}$, và $t = \frac{-v_0 + \sqrt{v_0^2 + 2a\Delta x}}{a}$.
      \end{flushleft}
   \end{tblr}
\end{table}

Khi mà chúng ta tính đến sự tái sử dụng của các kí hiệu, việc sử dụng các biểu thức khiến cho các khẳng định trở nên súc tích hơn nhiều.

Vậy thì giới hạn của khả năng sử dụng toán học trong khoa học nằm ở đâu? Nếu bài toán đang quan tâm là một bài toán định lượng, toán học là một công cụ hiệu quả để biểu diễn và giải nó, nhưng bài toán là định tính thì toán học lại trở nên không hữu dụng. Ví dụ, các công cụ giải tích có thể tính toán kích cỡ tối ưu khi sản xuất một hộp sữa để tiêu tốn ít giấy nhất, nhưng không thể giải thích tại sao bản thân nó có thể được sử dụng trong những bài toán như vậy. Phương trình toán học không có quá khứ, không có tương lai, không có cảm xúc hay cảm nhận vật lí. Góc nhìn của người làm toán, người làm khoa học mới đem lại câu chuyện cho các phương trình ấy.